\chapter{The Zeta Function}

\begin{center}\itshape
One would of course like to have a rigorous proof of this; but I have put aside the search for one, after some fleeting
vain attempts, as it is not necessary for the immediate purpose of my investigation.\\
\upshape --- Bernhard Riemann (1859), on the zeros of $\zeta$
\end{center}

The series $\zeta(s)=\sum_{n\ge1}n^{-s}$ has appeared in every chapter of this volume: its even values in Chapter~3, its
Bernoulli numbers in Chapter~1, its theta function in Chapter~10, its residues in Chapter~11. This chapter studies it
as a function of a complex variable. Euler discovered that it is a product over the primes, so that analysis of $\zeta$
is arithmetic in disguise. Riemann discovered, in an eight-page paper of 1859, that it extends to the whole complex plane
with a single pole and satisfies a symmetry $s\leftrightarrow1-s$, the \emph{functional equation}. We prove both, compute
$\zeta$ at the negative integers, explain what the famous ``$1+2+3+\cdots=-\frac1{12}$'' does and does not mean, and look at the
zeros on the line $\operatorname{Re}s=\frac12$, which Riemann conjectured to contain all the interesting ones.

What happens next (that $\zeta(1+it)\ne0$, and why this implies the prime number theorem) is the subject of Volume~6. Here we
keep a promise of Volume~2: Dirichlet series turn the convolution of arithmetic functions into multiplication.

\begin{think}
Expand $\bigl(1+\frac1{2^s}+\frac1{4^s}+\cdots\bigr)\bigl(1+\frac1{3^s}+\frac1{9^s}+\cdots\bigr)\bigl(1+\frac1{5^s}+\frac1{25^s}+\cdots\bigr)$. Which terms $\frac1{n^s}$ appear, and how often? What
must you multiply by to get every $\frac1{n^s}$ exactly once?
\end{think}

\section{Dirichlet series}

We write $s=\sigma+it$ with $\sigma,t$ real, and $n^{-s}=e^{-s\log n}$, so that $|n^{-s}|=n^{-\sigma}$. A \emph{Dirichlet series} is a series
$F(s)=\sum_{n\ge1}a_nn^{-s}$ with complex coefficients.

\begin{proposition}\label{prop:dirseries}
\begin{enumerate}[label=(\alph*)]
\item If $\sum|a_n|n^{-\sigma_0}<\infty$, then $\sum a_nn^{-s}$ converges absolutely and uniformly on $\operatorname{Re}s\ge\sigma_0$, and its sum is holomorphic on
$\operatorname{Re}s>\sigma_0$, with $F'(s)=-\sum a_n(\log n)n^{-s}$.
\item If the partial sums $A_N=\sum_{n\le N}a_n$ are bounded, then $\sum a_nn^{-s}$ converges on $\operatorname{Re}s>0$, uniformly on every closed disc there,
and its sum is holomorphic on $\operatorname{Re}s>0$.
\end{enumerate}
\end{proposition}
\begin{proof}
(a) $|a_nn^{-s}|\le|a_n|n^{-\sigma_0}$: apply the Weierstrass test, then Theorem~\ref{thm:holoparam}(a), each term being entire.

(b) For $x\ge n$, the difference $n^{-s}-x^{-s}=s\int_n^xu^{-s-1}\dd u$ has modulus at most $|s|(x-n)n^{-\sigma-1}$. By Proposition~\ref{prop:abel}, with $|A_n|\le M$,
\[
\sum_{n=1}^Na_nn^{-s}=A_NN^{-s}+\sum_{n=1}^{N-1}A_n\bigl(n^{-s}-(n+1)^{-s}\bigr).
\]
On a closed disc in $\operatorname{Re}s>0$ we have $|s|\le C$ and $\sigma\ge\sigma_1>0$, so $|A_NN^{-s}|\le MN^{-\sigma_1}\to0$ and the terms of the last sum are at most
$MCn^{-\sigma_1-1}$, summable. So the partial sums converge uniformly on the disc; apply Theorem~\ref{thm:holoparam}(a).
\end{proof}

For instance $\zeta(s)=\sum n^{-s}$ is holomorphic on $\operatorname{Re}s>1$, and the alternating series $\eta(s)=\sum(-1)^{n-1}n^{-s}$ is holomorphic on
$\operatorname{Re}s>0$ (Exercise~\ref{ex:eta}).

\begin{proposition}[products, uniqueness]\label{prop:dirmult}
\begin{enumerate}[label=(\alph*)]
\item If $\sum a_nn^{-s}$ and $\sum b_nn^{-s}$ converge absolutely, then their product is $\sum_nc_nn^{-s}$, with $c=a*b$ the Dirichlet convolution
$c_n=\sum_{de=n}a_db_e$, and this series converges absolutely.
\item If $\sum|a_n|n^{-\sigma_0}<\infty$ and $\sum a_nn^{-\sigma}=0$ for all real $\sigma>\sigma_0$, then $a_n=0$ for every $n$.
\end{enumerate}
\end{proposition}
\begin{proof}
(a) The double series $\sum_{d,e}a_db_e(de)^{-s}$ converges absolutely, being the product of two absolutely convergent series, so its terms
may be grouped by the value of $n=de$.

(b) Otherwise, let $m$ be the least index with $a_m\ne0$. For $\sigma>\sigma_0$, $0=m^\sigma\sum_na_nn^{-\sigma}=a_m+\sum_{n>m}a_n(m/n)^\sigma$, and
$|a_n(m/n)^\sigma|\le|a_n|n^{-\sigma_0}m^{\sigma_0}(m/n)^{\sigma-\sigma_0}$ for $n>m$. These terms tend to $0$ as $\sigma\to\infty$ and are dominated by the summable
$m^{\sigma_0}|a_n|n^{-\sigma_0}$, so the sum over $n>m$ tends to $0$ (dominated convergence for series). Hence $a_m=0$, a contradiction.
\end{proof}

So an absolutely convergent Dirichlet series determines its coefficients, and products of Dirichlet series compute Dirichlet
convolutions, exactly as products of power series compute Cauchy products. This is the promise of Volume~2, Chapter~4.

\section{The Euler product}

\begin{theorem}[Euler product]\label{thm:eulerproduct}
Let $f$ be a multiplicative arithmetic function with $\sum|f(n)|n^{-\sigma}<\infty$. Then for $\operatorname{Re}s\ge\sigma$,
\[
\sum_{n\ge1}\frac{f(n)}{n^s}=\prod_p\Bigl(1+\frac{f(p)}{p^s}+\frac{f(p^2)}{p^{2s}}+\cdots\Bigr),
\]
the product over the primes in increasing order being convergent. In particular, for $\operatorname{Re}s>1$,
\[
\zeta(s)=\prod_p\frac1{1-p^{-s}}.
\]
\end{theorem}
\begin{proof}
Each factor converges absolutely. For a bound $P$, multiplying the finitely many factors with $p\le P$ and expanding (absolutely
convergent series multiply term by term) gives $\sum f(p_1^{k_1})\cdots f(p_r^{k_r})(p_1^{k_1}\cdots p_r^{k_r})^{-s}$ over all choices of exponents. By unique
factorisation and multiplicativity, this is $\sum_{n\in S_P}f(n)n^{-s}$, where $S_P$ is the set of $n$ whose prime factors are all at most $P$. Every
$n\le P$ is in $S_P$, so the difference from the full series is at most $\sum_{n>P}|f(n)|n^{-\sigma}$, which tends to $0$. For $\zeta$, take $f=\mathbf 1$; the
factor is the geometric series $\sum_kp^{-ks}=\frac1{1-p^{-s}}$.
\end{proof}

The proof is the answer to the think box: the product over $p\le P$ produces each $n^{-s}$ once if $n$ has only prime factors
at most $P$, by unique factorisation. The Euler product is unique factorisation written analytically.

\begin{corollary}\label{cor:mobius}
For $\operatorname{Re}s>1$,
\[
\begin{gathered}\sum_{n\ge1}\frac{\mu(n)}{n^s}=\frac1{\zeta(s)},\qq\sum_{n\ge1}\frac{\varphi(n)}{n^s}=\frac{\zeta(s-1)}{\zeta(s)}\quad(\operatorname{Re}s>2).\end{gathered}
\]
In particular $\zeta(s)\ne0$ for $\operatorname{Re}s>1$.
\end{corollary}
\begin{proof}
By Proposition~\ref{prop:dirmult}(a), $\zeta(s)\sum\mu(n)n^{-s}=\sum(\mathbf 1*\mu)(n)n^{-s}=1$, since $\mu*\mathbf 1=\varepsilon$ (Volume~2, Theorem~6.4); so $\zeta(s)\ne0$.
Similarly $\varphi*\mathbf 1=\operatorname{id}$ (Volume~2, Proposition~4.5) gives $\zeta(s)\sum\varphi(n)n^{-s}=\sum n\cdot n^{-s}=\zeta(s-1)$, all series converging absolutely
for $\operatorname{Re}s>2$ because $\varphi(n)\le n$.
\end{proof}

At $s=2$: $\sum\frac{\mu(n)}{n^2}=\frac6{\pi^2}$. This is the value that Exercise~4.8 of Volume~2 needed: the probability that two random
integers up to $N$ are coprime tends to $\frac6{\pi^2}\approx0.6079$. The Euler product gives the heuristic behind it: two integers are
coprime when no prime divides both, and $p$ divides both with ``probability'' $\frac1{p^2}$, independently for different primes, so the
probability should be $\prod_p\bigl(1-\frac1{p^2}\bigr)=\frac1{\zeta(2)}$.

\begin{corollary}[Euler, 1737]\label{cor:primesdiverge}
$\displaystyle\sum_p\frac1p=\infty$.
\end{corollary}
\begin{proof}
Let $\sigma>1$ be real. By Theorem~\ref{thm:eulerproduct} and Proposition~\ref{prop:products}, $\log\zeta(\sigma)=\sum_p-\log(1-p^{-\sigma})$. For $0\le x\le\frac12$,
$-\log(1-x)=\sum_{k\ge1}\frac{x^k}k\le x+x^2\sum_{k\ge0}x^k\le x+2x^2$, so $\log\zeta(\sigma)\le\sum_pp^{-\sigma}+2\sum_nn^{-2}$. Since $\zeta(\sigma)\ge\int_1^\infty x^{-\sigma}\dd x=\frac1{\sigma-1}$,
the left side tends to $\infty$ as $\sigma\to1^+$, hence so does $\sum_pp^{-\sigma}$, which is at most $\sum_p\frac1p$.
\end{proof}

So the primes are much denser than the squares ($\sum\frac1{n^2}<\infty$). Exercise~\ref{ex:primesum} refines the argument to
$\sum_pp^{-\sigma}=\log\frac1{\sigma-1}+O(1)$, a first quantitative statement about primes from $\zeta$.

\section{Beyond the line of convergence}

The series diverges at $s=1$ and for real $s<1$. Yet $\zeta$ itself has a life beyond: it extends to a holomorphic function on a
larger set, and by the identity theorem (Theorem~\ref{thm:identity}) such an extension is unique. The first step is elementary.

\begin{proposition}\label{prop:zetastrip}
The function $\zeta(s)-\frac1{s-1}$ extends to a holomorphic function on $\operatorname{Re}s>0$, whose value at $s=1$ is Euler's constant $\gamma$. So $\zeta$ is
holomorphic on $\operatorname{Re}s>0$ except for a simple pole at $1$ with residue $1$, and $\zeta(s)=\frac1{s-1}+\gamma+O(s-1)$ near $1$.
\end{proposition}
\begin{proof}
For $\operatorname{Re}s>1$, $\frac1{s-1}=\int_1^\infty x^{-s}\dd x=\sum_n\int_n^{n+1}x^{-s}\dd x$, so
\[
\begin{gathered}\zeta(s)-\frac1{s-1}=\sum_{n\ge1}\phi_n(s),\qq\phi_n(s)=\int_n^{n+1}\bigl(n^{-s}-x^{-s}\bigr)\dd x.\end{gathered}
\]
Each $\phi_n$ is entire (Theorem~\ref{thm:holoparam}(b)), and by the bound in the proof of Proposition~\ref{prop:dirseries}(b), $|\phi_n(s)|\le|s|n^{-\sigma-1}$. So the series
converges uniformly on every closed disc in $\operatorname{Re}s>0$ and defines a holomorphic function there (Theorem~\ref{thm:holoparam}(a)). At $s=1$,
$\sum_{n\le N}\phi_n(1)=H_N-\log(N+1)$, which tends to $\gamma$ (Volume~1, Proposition~3.16).
\end{proof}

\begin{figure}[ht]
\centering
\fig{ch10-plane}
\caption{The $s$-plane. On $\operatorname{Re}s>1$ (charcoal) the series and the Euler product converge absolutely and $\zeta$ has no zeros. The pole is at $1$.
The trivial zeros (ochre) are at $-2,-4,-6,\dots$ All other zeros lie in the critical strip $0\le\operatorname{Re}s\le1$ (coral), symmetric about the critical
line $\operatorname{Re}s=\frac12$ (dashed); the first ones, in rust, are at $\frac12\pm14.13i$, $\frac12\pm21.02i$, $\frac12\pm25.01i$.}
\label{fig:plane}
\end{figure}

\section{Riemann's functional equation}

Riemann's idea is to write $n^{-s}$ as an integral of the Gaussian $e^{-\pi n^2t}$ and to sum over $n$, which brings in the theta function of
Chapter~10 and its transformation law.

\begin{lemma}\label{lem:mellingauss}
For $\operatorname{Re}s>0$ and $n\ge1$, $\displaystyle\int_0^\infty e^{-\pi n^2t}t^{s/2-1}\dd t=\pi^{-s/2}\,\Gamma\Bigl(\frac s2\Bigr)\,n^{-s}$.
\end{lemma}
\begin{proof}
Substitute $u=\pi n^2t$: $t^{s/2-1}=(\pi n^2)^{1-s/2}u^{s/2-1}$ (powers of positive reals multiply), and the integral becomes $(\pi n^2)^{-s/2}\Gamma(\frac s2)$ by
Theorem~\ref{thm:gammac}(a).
\end{proof}

Let $\omega(t)=\sum_{n\ge1}e^{-\pi n^2t}=\frac{\theta(t)-1}2$ for $t>0$, where $\theta$ is the theta function of Corollary~\ref{cor:theta}. For $t\ge1$, $\omega(t)\le\sum_{n\ge1}e^{-\pi nt}\le2e^{-\pi t}$,
since $e^{-\pi}<\frac12$. For $0<t\le1$, the theta transformation gives $\omega(t)\le\frac12\theta(t)=\frac12t^{-1/2}\theta(1/t)\le\frac12\theta(1)t^{-1/2}$.

\begin{theorem}[Riemann, 1859]\label{thm:functional}
Let $\Lambda(s)=\pi^{-s/2}\Gamma(\frac s2)\zeta(s)$ for $\operatorname{Re}s>1$. Then
\[
\begin{aligned}\Lambda(s)&=-\frac1s-\frac1{1-s}\eqbr\quad+\int_1^\infty\omega(t)\bigl(t^{s/2}+t^{(1-s)/2}\bigr)\frac{\mathrm dt}t.\end{aligned}
\]
The integral is an entire function of $s$, so $\Lambda$ extends to a holomorphic function on $\C\setminus\{0,1\}$, with simple poles at $0$ and $1$, and
\[
\Lambda(1-s)=\Lambda(s).
\]
\end{theorem}
\begin{proof}
Let $\operatorname{Re}s=\sigma>1$. By Lemma~\ref{lem:mellingauss}, $\Lambda(s)=\sum_n\int_0^\infty e^{-\pi n^2t}t^{s/2-1}\dd t$. The partial sums of $\sum_ne^{-\pi n^2t}t^{s/2-1}$ converge to $\omega(t)t^{s/2-1}$
uniformly on every closed bounded subinterval of $\left]0,\infty\right[$ and are dominated by $\omega(t)t^{\sigma/2-1}$, which is integrable: it is at most $\frac12\theta(1)t^{\sigma/2-3/2}$ near $0$
(integrable since $\sigma>1$) and at most $2e^{-\pi t}t^{\sigma/2-1}$ near $\infty$. By Corollary~\ref{cor:dct}, $\Lambda(s)=\int_0^\infty\omega(t)t^{s/2-1}\dd t$.

Split the integral at $1$, and substitute $t=\frac1u$ in the part $J$ over $\left]0,1\right]$. The theta transformation $\theta(\frac1u)=\sqrt u\,\theta(u)$ reads
$\omega(\frac1u)=\sqrt u\,\omega(u)+\frac{\sqrt u}2-\frac12$, so
\[
\begin{aligned}J&=\int_1^\infty\omega\Bigl(\frac1u\Bigr)u^{-s/2-1}\dd u\\&=\int_1^\infty\omega(u)u^{(1-s)/2-1}\dd u\\&\quad+\frac12\int_1^\infty\bigl(u^{-(s+1)/2}-u^{-s/2-1}\bigr)\dd u.\end{aligned}
\]
The integral without the factor $\frac12$ equals $\frac2{s-1}-\frac2s$ (the antiderivative of $u^{-(s+1)/2}$ contributes $\frac2{s-1}$, and that of $u^{-s/2-1}$ contributes $\frac2s$). With the factor already written, the term is $\frac12\bigl(\frac2{s-1}-\frac2s\bigr)=\frac1{s-1}-\frac1s=-\frac1s-\frac1{1-s}$. Adding the part over $[1,\infty[$ gives the formula.

The integral $I(s)=\int_1^\infty\omega(t)(t^{s/2}+t^{(1-s)/2})\frac{\mathrm dt}t$ is entire by Theorem~\ref{thm:holoparam}(b): on a closed disc where $|\operatorname{Re}s|\le K$, the integrand is
dominated by $4e^{-\pi t}t^{K/2+1/2}$. The right side of the formula is therefore holomorphic on $\C\setminus\{0,1\}$ and equals $\Lambda$ on $\operatorname{Re}s>1$; this defines
the extension. Replacing $s$ by $1-s$ exchanges $-\frac1s$ and $-\frac1{1-s}$ and leaves $I$ unchanged.
\end{proof}

\begin{beauty}
The functional equation is the theta transformation in disguise, and the theta transformation is Poisson summation for
the Gaussian, which is its own Fourier transform. So the symmetry $s\leftrightarrow1-s$ of $\zeta$ comes from the self-duality of the
integers under the Fourier transform (Chapter~10). Riemann's 1859 paper, which also introduced the zeros and their link
with the primes, contains this argument almost word for word.
\end{beauty}

\begin{corollary}\label{cor:zetaplane}
\begin{enumerate}[label=(\alph*)]
\item $\zeta$ extends to a function holomorphic on $\C\setminus\{1\}$, with a simple pole of residue $1$ at $1$. Moreover $\zeta(\overline s)=\overline{\zeta(s)}$.
\item $\zeta(0)=-\frac12$, and $\zeta(-2k)=0$ for every $k\ge1$: these are simple zeros, the \emph{trivial zeros}.
\item For $\operatorname{Re}s<0$, $\zeta(s)=0$ only at the trivial zeros.
\item (asymmetric form) For every $s$ that is not an integer $\le1$,
\[
\zeta(1-s)=2(2\pi)^{-s}\cos\Bigl(\frac{\pi s}2\Bigr)\Gamma(s)\zeta(s).
\]
The \emph{completed} function $\xi(s)=\frac12s(s-1)\Lambda(s)$ is entire, with $\xi(1-s)=\xi(s)$.
\end{enumerate}
\end{corollary}
\begin{proof}
(a) $\zeta(s)=\pi^{s/2}\frac1{\Gamma(s/2)}\Lambda(s)$, where $\frac1\Gamma$ is entire (Theorem~\ref{thm:gammac}(c)), so $\zeta$ is holomorphic except possibly at $0$ and $1$. At $1$,
$\pi^{1/2}/\Gamma(\frac12)=1$ and the residue is that of $\Lambda$, namely $1$ (Proposition~\ref{prop:zetastrip}). Conjugation: $s\mapsto\overline{\zeta(\overline s)}$ is holomorphic and
equals $\zeta$ on $\left]1,\infty\right[$; apply Theorem~\ref{thm:identity}.

(b) Near $0$, $\Lambda(s)=-\frac1s+O(1)$ and $\frac1{\Gamma(s/2)}=\frac s2+O(s^2)$ (since $\Gamma(z)=\frac1z+O(1)$ by Theorem~\ref{thm:gammac}(b)); so $\zeta(0)=-\frac12$. At $s=-2k$, $\Lambda$ is
holomorphic and $\frac1{\Gamma(s/2)}$ has a simple zero, so $\zeta(-2k)=0$. The zero is simple because $\Lambda(-2k)=\Lambda(1+2k)=\pi^{-k-1/2}\Gamma(k+\frac12)\zeta(2k+1)\ne0$.

(c) If $\operatorname{Re}s<0$, then $\operatorname{Re}(1-s)>1$, so $\Lambda(s)=\Lambda(1-s)\ne0$ by Corollary~\ref{cor:mobius} (and $\Gamma$ has no zeros). Hence $\zeta(s)=0$ only where $\frac1{\Gamma(s/2)}=0$.

(d) Let $R(s)=\frac{\Gamma(s/2)}{\Gamma((1-s)/2)}$. By the functional equation, $\zeta(1-s)=\pi^{1/2-s}R(s)\zeta(s)$. The reflection formula gives $\Gamma(\frac{1-s}2)\Gamma(\frac{1+s}2)=\frac\pi{\cos(\pi s/2)}$, and
the duplication formula (Theorem~\ref{thm:duplication}, extended to $\C$ by the identity theorem) gives $\Gamma(\frac s2)\Gamma(\frac{s+1}2)=2^{1-s}\sqrt\pi\,\Gamma(s)$. So
$R(s)=\frac{2^{1-s}\sqrt\pi\,\Gamma(s)\cos(\pi s/2)}\pi$, which gives the formula. Finally $\frac12s(s-1)$ cancels the poles of $\Lambda$ and is invariant under
$s\mapsto1-s$.
\end{proof}

\begin{theorem}\label{thm:zetaneg}
For every $k\ge1$, $\displaystyle\zeta(1-2k)=-\frac{B_{2k}}{2k}$. Thus $\zeta(-1)=-\frac1{12}$, $\zeta(-3)=\frac1{120}$, $\zeta(-5)=-\frac1{252}$, $\zeta(-7)=\frac1{240}$.
\end{theorem}
\begin{proof}
Put $s=2k$ in Corollary~\ref{cor:zetaplane}(d), where $\cos(\pi k)=(-1)^k$ and $\Gamma(2k)=(2k-1)!$, and use Euler's formula (Theorem~\ref{thm:zeta2k}):
\[
\begin{aligned}\zeta(1-2k)&=\frac{2(-1)^k(2k-1)!}{(2\pi)^{2k}}\eqbr\quad\cdot\frac{(-1)^{k+1}B_{2k}(2\pi)^{2k}}{2(2k)!}\eqbr=-\frac{(2k-1)!}{(2k)!}B_{2k}=-\frac{B_{2k}}{2k}.\end{aligned}
\]
The values follow from $B_2=\frac16$, $B_4=-\frac1{30}$, $B_6=\frac1{42}$, $B_8=-\frac1{30}$.
\end{proof}

With $\zeta(0)=-\frac12=-B_1(1)$ and $\zeta(-2k)=0=-\frac{B_{2k+1}}{2k+1}$, the single formula $\zeta(-n)=-\frac{B_{n+1}(1)}{n+1}$ covers all $n\ge0$; Exercise~\ref{ex:hurwitz}
proves it directly, and its generalisation to the Hurwitz zeta function, without the functional equation. Euler had found these
values in 1749 by summing divergent series, and had guessed the functional equation from them.

\section{What $1+2+3+\cdots=-\frac1{12}$ means}

The series $\sum n$ diverges, and nothing will change that. The equation $\zeta(-1)=-\frac1{12}$ says something else: the analytic continuation of
$\sum n^{-s}$ takes the value $-\frac1{12}$ at $s=-1$. There is a concrete way to see this number appear from the divergent series itself:
smooth the cut-off.

\begin{proposition}\label{prop:smoothed}
As $N\to\infty$,
\[
\begin{gathered}\sum_{n\ge1}e^{-n/N}=N-\frac12+O\Bigl(\frac1N\Bigr),\qq\sum_{n\ge1}ne^{-n/N}=N^2-\frac1{12}+O\Bigl(\frac1{N^2}\Bigr).\end{gathered}
\]
\end{proposition}
\begin{proof}
Let $h=\frac1N$. The first sum is $\frac1{e^h-1}=\frac1h\cdot\frac h{e^h-1}=\frac1h\bigl(1-\frac h2+O(h^2)\bigr)$. The second is $\sum nx^n=\frac x{(1-x)^2}$ at $x=e^{-h}$, that is
$\frac{e^{-h}}{(1-e^{-h})^2}=\frac1{(e^{h/2}-e^{-h/2})^2}=\frac1{4\sinh^2(h/2)}$. Since $\sinh u=u\bigl(1+\frac{u^2}6+O(u^4)\bigr)$, $4\sinh^2\frac h2=h^2\bigl(1+\frac{h^2}{12}+O(h^4)\bigr)$, whose inverse is
$\frac1{h^2}-\frac1{12}+O(h^2)$.
\end{proof}

The divergent part $N^2$ depends on how we cut off the sum. For this exponential cut-off the constant term is $-\frac1{12}=\zeta(-1)$, the value of the analytic continuation. The same finite part for a general smooth cut-off, and for $\sum n^k$, is a theorem of the Mellin transform; it is not proved here. With a sharp cut-off, $\sum_{n\le N}n=\frac{N^2}2+\frac N2$, there is no constant
term at all: the sharp cut-off is too rough for the finite part to show. Physicists use exactly this finite part when they
compute the Casimir force between two plates.

\begin{think}
Is $1-1+1-1+\cdots=\frac12$? Compute $\sum_{n\ge0}(-1)^ne^{-n/N}$. Then compare with $\eta(0)$, where $\eta$ is the alternating zeta function of
Exercise~\ref{ex:eta}.
\end{think}

\section{The zeros}

By Corollaries~\ref{cor:mobius} and~\ref{cor:zetaplane}, all zeros of $\zeta$ other than the trivial ones lie in the \emph{critical strip} $0\le\operatorname{Re}s\le1$,
symmetrically with respect to the real axis ($\zeta(\overline s)=\overline{\zeta(s)}$) and to the \emph{critical line} $\operatorname{Re}s=\frac12$ (the functional equation, as $\Gamma$
has no zeros). Exercise~\ref{ex:eta} shows that there are none on the real segment $\left]0,1\right[$.

On the critical line, the functional equation gives $\Lambda(\frac12+it)=\Lambda(\frac12-it)=\overline{\Lambda(\frac12+it)}$: \emph{$\Lambda$ is real on the critical line}. Writing
$\Gamma(\frac14+\frac{it}2)=|\Gamma(\frac14+\frac{it}2)|e^{i\vartheta_0(t)}$ with $\vartheta_0$ continuous and $\vartheta(t)=\vartheta_0(t)-\frac t2\log\pi$, we get
\[
\begin{gathered}\Lambda\Bigl(\frac12+it\Bigr)=\pi^{-1/4}\Bigl|\Gamma\Bigl(\frac14+\frac{it}2\Bigr)\Bigr|\,Z(t),\qq Z(t)=e^{i\vartheta(t)}\zeta\Bigl(\frac12+it\Bigr).\end{gathered}
\]
So \emph{Hardy's function} $Z$ is real-valued, with $|Z(t)|=|\zeta(\frac12+it)|$, and every sign change of $Z$ is a zero of $\zeta$ on the critical
line (Figure~\ref{fig:hardyz}). This is how zeros are found and certified on a computer: by the intermediate value theorem.

\begin{figure}[ht]
\centering
\fig{ch10-hardyz}
\caption{Hardy's function $Z(t)$ (charcoal), real-valued with $|Z(t)|=|\zeta(\frac12+it)|$ (grey), for $0\le t\le50$. Its sign changes are zeros of $\zeta$ on the
critical line: the ten zeros with $0<t<50$, at $t=14.13$, $21.02$, $25.01$, $30.42$, $32.94$, $37.59$, $40.92$, $43.33$, $48.01$, $49.77$ (rust).}
\label{fig:hardyz}
\end{figure}

\begin{remark}[the Riemann hypothesis]
Riemann conjectured that all non-trivial zeros lie on the critical line. Hardy proved in 1914 that infinitely many do, and
Selberg in 1942 that a positive proportion do. The first $10^{13}$ zeros have been computed and lie on the line, and the
hypothesis has been verified for all zeros with $0<t\le3\cdot10^{12}$ (Platt and Trudgian, 2021). It remains open, and is one of
the Millennium Prize Problems. Its arithmetic meaning, and the proof that there are no zeros on the line $\operatorname{Re}s=1$ (which gives the
prime number theorem), are the subject of Volume~6.
\end{remark}

\section{Exercises}

\begin{exercise}[$\star$]\label{ex:tau}
Show that for $\operatorname{Re}s>2$, $\sum\frac{\tau(n)}{n^s}=\zeta(s)^2$, $\sum\frac{\sigma(n)}{n^s}=\zeta(s)\zeta(s-1)$ and $\sum\frac{|\mu(n)|}{n^s}=\frac{\zeta(s)}{\zeta(2s)}$.
\end{exercise}

\begin{exercise}[$\star\star$]\label{ex:squarefree}
Let $Q(N)$ be the number of squarefree integers in $\{1,\dots,N\}$. Show that $Q(N)=\sum_{d\le\sqrt N}\mu(d)\lfloor N/d^2\rfloor$ and that $\frac{Q(N)}N\to\frac6{\pi^2}$.
\end{exercise}

\begin{exercise}[$\star\star$]\label{ex:eta}
Let $\eta(s)=\sum_{n\ge1}(-1)^{n-1}n^{-s}$.
\begin{enumerate}[label=(\alph*)]
\item Show that $\eta$ is holomorphic on $\operatorname{Re}s>0$ and that $\eta(s)=(1-2^{1-s})\zeta(s)$ there.
\item Deduce that $\zeta(\sigma)<0$ for $0<\sigma<1$, and recover the residue of $\zeta$ at $1$ from $\eta(1)=\log2$.
\end{enumerate}
\end{exercise}

\begin{exercise}[$\star\star$]\label{ex:primesum}
Show that $\log\zeta(\sigma)=\log\frac1{\sigma-1}+O(\sigma-1)$ as $\sigma\to1^+$, and that $\sum_pp^{-\sigma}=\log\frac1{\sigma-1}+O(1)$.
\end{exercise}

\begin{exercise}[$\star\star\star$]\label{ex:hurwitz}
(The Hurwitz zeta function.) For $0<a\le1$ and $\operatorname{Re}s>1$ let $\zeta(s,a)=\sum_{n\ge0}(n+a)^{-s}$, so that $\zeta(s,1)=\zeta(s)$, and let $g(t)=\frac{e^{-at}}{1-e^{-t}}$.
\begin{enumerate}[label=(\alph*)]
\item Show that $\Gamma(s)\zeta(s,a)=\int_0^\infty t^{s-1}g(t)\dd t$.
\item Show that $tg(t)=\sum_{m\ge0}(-1)^mB_m(a)\frac{t^m}{m!}$ for $|t|<2\pi$.
\item Deduce that $\zeta(s,a)$ extends to a function holomorphic on $\C\setminus\{1\}$, and that $\zeta(-n,a)=-\frac{B_{n+1}(a)}{n+1}$ for every integer $n\ge0$.
\end{enumerate}
\end{exercise}

\begin{exercise}[$\star\star$]\label{ex:zeta2fe}
Using only Exercise~\ref{ex:hurwitz} with $a=1$ and the asymmetric functional equation, find $\zeta(2)$ and $\zeta(4)$. Explain why this gives a new proof
of Theorem~\ref{thm:zeta2k}.
\end{exercise}

\begin{exercise}[Journey]\label{pb:gausskuzmin}
\emph{The Gauss--Kuzmin theorem.} Let $x$ be uniform in $]0,1[$, $m_n(t)=P(T^nx<t)$ and $f_n=m_n'$, so $f_0=1$.
\begin{enumerate}[label=(\alph*)]
\item Show that $m_{n+1}(t)=\sum_{k\ge1}\bigl(m_n(\frac1k)-m_n(\frac1{k+t})\bigr)$, and deduce
$f_{n+1}(t)=\sum_{k\ge1}\frac1{(k+t)^2}f_n\bigl(\frac1{k+t}\bigr)$, with every $f_n$ continuously differentiable.
\item Let $\phi_n(t)=(1+t)f_n(t)$. Show that $\phi_{n+1}(t)=\sum_kP_k(t)\phi_n(u_k)$ with $u_k=\frac1{k+t}$ and
$P_k(t)=\frac{1+t}{(k+t)(k+t+1)}$, and that $\sum_kP_k(t)=1$.
\item Show that $\sup|\phi_{n+1}'|\le c\,\sup|\phi_n'|$ with
\[
c=\frac12+2\sum_{k\ge2}\frac1{k^3(k+1)}+\sum_{k\ge2}\frac{k-1}{k(k+1)^2}=2\zeta(3)-\frac32=0.9041\ldots
\]
\item Deduce that $\Bigl|f_n(t)-\frac1{(1+t)\log2}\Bigr|\le\frac{c^n}{1+t}$, and conclude: $|m_n(t)-\log_2(1+t)|\le c^n\log(1+t)$.
\item Deduce that $P(a_n=k)\to\log_2\frac{(k+1)^2}{k(k+2)}$ for $x$ uniform.
\end{enumerate}
\end{exercise}

\begin{exercise}[$\star\star$]\label{pb:sqrtsum}
\emph{The constant in $1+\sqrt2+\dots+\sqrt n$.} Use the continuation of $\zeta$ by Euler--Maclaurin (Volume~5, Proposition~13.7) to show that
\[
\begin{aligned}\sum_{k=1}^n\sqrt k&=\frac23n^{3/2}+\frac12\sqrt n+\zeta\Bigl(-\frac12\Bigr)+\frac1{24\sqrt n}+O\bigl(n^{-5/2}\bigr),\\
\sum_{k=1}^n\frac1{\sqrt k}&=2\sqrt n+\zeta\Bigl(\frac12\Bigr)+\frac1{2\sqrt n}-\frac1{24n^{3/2}}+O\bigl(n^{-7/2}\bigr),\end{aligned}
\]
where $\zeta(-\frac12)=-0.2078862\ldots$ and $\zeta(\frac12)=-1.4603545\ldots$ So the constant left over when the growing terms are removed is a value of the zeta function.
\end{exercise}

\section{Hints}

\begin{hint}{ex:tau}
$\tau=\mathbf 1*\mathbf 1$, $\sigma=\operatorname{id}*\mathbf 1$; for the last one, use the Euler product.
\end{hint}
\begin{hint}{ex:squarefree}
$\sum_{d^2\mid n}\mu(d)$ is $1$ if $n$ is squarefree and $0$ otherwise. Bound the error made by dropping the floor.
\end{hint}
\begin{hint}{ex:eta}
(a) Proposition~\ref{prop:dirseries}(b); for the identity, $\zeta(s)-\eta(s)=2\sum_n(2n)^{-s}$. (b) An alternating series with decreasing terms has the sign of its
first term.
\end{hint}
\begin{hint}{ex:primesum}
$\frac1{\sigma-1}\le\zeta(\sigma)\le1+\frac1{\sigma-1}$. Then the inequality $x\le-\log(1-x)\le x+2x^2$ from Corollary~\ref{cor:primesdiverge}.
\end{hint}
\begin{hint}{ex:hurwitz}
(a) As in Theorem~\ref{thm:functional}. (b) $tg(t)=\frac{te^{(1-a)t}}{e^t-1}$ and the symmetry $B_m(1-a)=(-1)^mB_m(a)$, read off the generating function. (c) Subtract the first $M$
terms of the expansion on $\left]0,1\right]$, keep $\int_1^\infty$ as it is, and use $\frac1{\Gamma(s)}=(-1)^nn!\,(s+n)+O((s+n)^2)$ near $-n$.
\end{hint}
\begin{hint}{ex:zeta2fe}
Put $s=2$ and $s=4$ in Corollary~\ref{cor:zetaplane}(d). The proof of that corollary does not use Theorem~\ref{thm:zeta2k}.
\end{hint}

\begin{hint}{pb:gausskuzmin}
(c) Differentiate the recursion of (b). Since $\sum P_k'=0$, the first sum is $\sum P_k'(t)\bigl(\phi_n(u_k)-\phi_n(u_1)\bigr)$.
(d) $\frac1{\log2}$ is a weighted average of $\phi_n$.
\end{hint}

\begin{hint}{pb:sqrtsum}
In Proposition~11.7 of Volume~5, take $s=-\frac12$ (respectively $s=\frac12$) and $M=2$, and replace $N$ by $n$. Bound the integral remainder.
\end{hint}

\section{Solutions}

\begin{solution}{ex:tau}
\idea Write each function as a convolution of functions whose Dirichlet series are known, and multiply.
\writeup All series converge absolutely for $\operatorname{Re}s>2$, since $\tau(n)\le n$ and $\sigma(n)\le n^2$. By Proposition~\ref{prop:dirmult}(a), $\tau=\mathbf 1*\mathbf 1$ gives $\zeta(s)^2$, and
$\sigma=\operatorname{id}*\mathbf 1$ gives $\zeta(s-1)\zeta(s)$, since $\sum n\cdot n^{-s}=\zeta(s-1)$. For $|\mu|$, the Euler product of Theorem~\ref{thm:eulerproduct} has factors $1+p^{-s}=\frac{1-p^{-2s}}{1-p^{-s}}$,
and the product of the numerators is $\frac1{\zeta(2s)}$, that of the denominators' inverses $\zeta(s)$. (This one holds for $\operatorname{Re}s>1$.)
\end{solution}

\begin{solution}{ex:squarefree}
\idea Detect squarefree numbers with M\"obius, exactly as Volume~2 detected coprime pairs, and let Corollary~\ref{cor:mobius} supply the constant.
\writeup Write $n=m^2k$ with $k$ squarefree; then $d^2\mid n$ iff $d\mid m$, so $\sum_{d^2\mid n}\mu(d)=\sum_{d\mid m}\mu(d)=[m=1]$ by Volume~2, Theorem~6.4, and $m=1$ means $n$ is squarefree.
Summing over $n\le N$ and exchanging the sums, $Q(N)=\sum_{d\le\sqrt N}\mu(d)\lfloor N/d^2\rfloor$. Replacing $\lfloor N/d^2\rfloor$ by $N/d^2$ costs at most $1$ per term, so
\[
\Bigl|Q(N)-N\sum_{d\le\sqrt N}\frac{\mu(d)}{d^2}\Bigr|\le\sqrt N.
\]
Dividing by $N$, $\frac{Q(N)}N\to\sum_d\frac{\mu(d)}{d^2}=\frac1{\zeta(2)}=\frac6{\pi^2}$. So about $61\%$ of integers are squarefree, the same proportion as coprime pairs.
\end{solution}

\begin{solution}{ex:eta}
\idea The alternating signs give bounded partial sums, which is all Proposition~\ref{prop:dirseries}(b) needs; and removing the even terms twice turns
$\zeta$ into $\eta$.
\writeup (a) The partial sums of $(-1)^{n-1}$ are $0$ or $1$, so $\eta$ is holomorphic on $\operatorname{Re}s>0$. For $\operatorname{Re}s>1$, $\zeta(s)-\eta(s)=2\sum_n(2n)^{-s}=2^{1-s}\zeta(s)$, so
$\eta(s)=(1-2^{1-s})\zeta(s)$. Both sides are holomorphic on $\operatorname{Re}s>0$ minus $\{1\}$ (Proposition~\ref{prop:zetastrip}), a connected set, so the identity holds there.

(b) For real $0<\sigma<1$, $\eta(\sigma)=(1-2^{-\sigma})+(3^{-\sigma}-4^{-\sigma})+\cdots>0$, a series of positive terms. And $1-2^{1-\sigma}<0$. So $\zeta(\sigma)=\frac{\eta(\sigma)}{1-2^{1-\sigma}}<0$; in
particular $\zeta(\frac12)\approx-1.4604$. Near $1$, $1-2^{1-s}=(s-1)\log2+O((s-1)^2)$, so $(s-1)\zeta(s)=\frac{(s-1)\eta(s)}{1-2^{1-s}}\to\frac{\eta(1)}{\log2}=1$, using
$\eta(1)=1-\frac12+\frac13-\cdots=\log2$ (Volume~1) and the continuity of $\eta$ at $1$.
\end{solution}

\begin{solution}{ex:primesum}
\idea $\zeta(\sigma)$ is squeezed between two integrals, and its logarithm is $\sum_pp^{-\sigma}$ up to a bounded error.
\writeup Comparing the sum with integrals of the decreasing function $x^{-\sigma}$, $\int_1^\infty x^{-\sigma}\dd x\le\zeta(\sigma)\le1+\int_1^\infty x^{-\sigma}\dd x$, that is
$\frac1{\sigma-1}\le\zeta(\sigma)\le\frac\sigma{\sigma-1}$. Taking logarithms, $0\le\log\zeta(\sigma)-\log\frac1{\sigma-1}\le\log\sigma\le\sigma-1$. Next, by Corollary~\ref{cor:primesdiverge}'s inequalities,
$\sum_pp^{-\sigma}\le\log\zeta(\sigma)\le\sum_pp^{-\sigma}+2\zeta(2)$. Combining, $\bigl|\sum_pp^{-\sigma}-\log\frac1{\sigma-1}\bigr|\le2\zeta(2)+1$ for $1<\sigma\le2$. Volume~6 turns this into Mertens's
theorem $\sum_{p\le x}\frac1p=\log\log x+O(1)$.
\end{solution}

\begin{solution}{ex:hurwitz}
\idea Riemann's second method: the Mellin transform turns the sum into an integral of $t^{s-1}$ against a function whose Taylor
expansion at $0$ has Bernoulli coefficients. Each term of that expansion produces one pole of $\Gamma(s)\zeta(s,a)$, and dividing by $\Gamma$
leaves a finite value at each negative integer, read off from the Bernoulli polynomials.
\writeup (a) As in Lemma~\ref{lem:mellingauss}, $\Gamma(s)(n+a)^{-s}=\int_0^\infty t^{s-1}e^{-(n+a)t}\dd t$. Summing over $n$, the partial sums of $\sum_nt^{s-1}e^{-(n+a)t}$ tend to $t^{s-1}g(t)$
uniformly on closed bounded subintervals of $\left]0,\infty\right[$ and are dominated by $t^{\sigma-1}g(t)$, which is integrable for $\sigma>1$ since $g(t)\le\frac{e^{-at}}t\cdot\frac t{1-e^{-t}}=O(t^{-1})$
at $0$ and $g(t)=O(e^{-at})$ at $\infty$. Corollary~\ref{cor:dct} gives the formula.

(b) $tg(t)=\frac{te^{(1-a)t}}{e^t-1}=\sum B_m(1-a)\frac{t^m}{m!}$ by the definition of the Bernoulli polynomials (Chapter~1), and this power series converges for $|t|<2\pi$
(Corollary~\ref{cor:radius}, the only singularities being at $2\pi i\Z\setminus\{0\}$). Replacing $t$ by $-t$ in $\frac{te^{xt}}{e^t-1}$ with $x=1-a$ gives $\frac{-te^{-(1-a)t}}{e^{-t}-1}=\frac{te^{at}}{e^t-1}$,
so $B_m(1-a)(-1)^m=B_m(a)$.

(c) Let $c_m=\frac{(-1)^mB_m(a)}{m!}$ and fix $M\ge1$. Then $r_M(t)=g(t)-\sum_{m=0}^Mc_mt^{m-1}$ is $O(t^M)$ on $\left]0,1\right]$ by (b). For $\operatorname{Re}s>1$,
\[
\begin{aligned}\Gamma(s)\zeta(s,a)&=\sum_{m=0}^M\frac{c_m}{s+m-1}+\int_0^1t^{s-1}r_M(t)\dd t\eqbr\quad+\int_1^\infty t^{s-1}g(t)\dd t,\end{aligned}
\]
using $\int_0^1t^{s+m-2}\dd t=\frac1{s+m-1}$. The last integral is entire and the middle one is holomorphic for $\operatorname{Re}s>-M$, by Theorem~\ref{thm:holoparam}(b). So the right
side is holomorphic on $\operatorname{Re}s>-M$ except for simple poles at $s=1-m$, $0\le m\le M$. Multiplying by the entire $\frac1\Gamma$, which vanishes at $0,-1,-2,\dots$, the poles at
$s=1-m$ with $m\ge1$ disappear; only $s=1$ remains. As $M$ is arbitrary, $\zeta(s,a)$ is holomorphic on $\C\setminus\{1\}$. At $s=-n$ the only singular term is $\frac{c_{n+1}}{s+n}$,
and $\frac1{\Gamma(s)}=(-1)^nn!\,(s+n)+O((s+n)^2)$ since $\Gamma$ has residue $\frac{(-1)^n}{n!}$ there (Theorem~\ref{thm:gammac}(b)). Hence
\[
\begin{aligned}\zeta(-n,a)&=(-1)^nn!\,c_{n+1}\eqbr=(-1)^nn!\frac{(-1)^{n+1}B_{n+1}(a)}{(n+1)!}\eqbr=-\frac{B_{n+1}(a)}{n+1}.\end{aligned}
\]
For $a=1$: $\zeta(-n)=-\frac{B_{n+1}(1)}{n+1}$, which is $-\frac12$ for $n=0$ ($B_1(1)=\frac12$), and $-\frac{B_{n+1}}{n+1}$ for $n\ge1$ by Proposition~\ref{prop:bernoulli}(c).
\end{solution}

\begin{solution}{ex:zeta2fe}
\idea The functional equation links $\zeta(2k)$ to $\zeta(1-2k)$, and the Hurwitz computation gives $\zeta(1-2k)$ from Bernoulli numbers without knowing
$\zeta(2k)$. So the functional equation transports Bernoulli numbers from the left half-plane to the right.
\writeup At $s=2$, since $\cos\pi=-1$ and $\Gamma(2)=1$, the equation reads $\zeta(-1)=-\frac{2\zeta(2)}{(2\pi)^2}=-\frac{\zeta(2)}{2\pi^2}$; and $\zeta(-1)=-\frac{B_2}2=-\frac1{12}$, so $\zeta(2)=\frac{\pi^2}6$. At $s=4$,
with $\cos2\pi=1$ and $\Gamma(4)=6$, $\zeta(-3)=\frac{12\,\zeta(4)}{(2\pi)^4}=\frac{3\zeta(4)}{4\pi^4}$, and $\zeta(-3)=-\frac{B_4}4=\frac1{120}$, so $\zeta(4)=\frac{\pi^4}{90}$. In general, $s=2k$ gives
$-\frac{B_{2k}}{2k}=\frac{2(-1)^k(2k-1)!}{(2\pi)^{2k}}\zeta(2k)$, which is Euler's formula. The proof of Corollary~\ref{cor:zetaplane}(d) used the theta function, $\Gamma$, reflection and duplication,
never Theorem~\ref{thm:zeta2k}; and Exercise~\ref{ex:hurwitz} used only the Bernoulli generating function. So this is an independent proof of Euler's formula.
\end{solution}

\begin{solution}{pb:gausskuzmin}
\idea The transfer operator replaces $\phi_n$ by an average of its values at the points $u_k$, and these points are squeezed
together: $|u_k'|=(k+t)^{-2}$. Averaging over squeezed points contracts the Lipschitz constant, and a function with tiny
Lipschitz constant is almost constant; the constant is forced by normalisation.

\writeup (a) For irrational $x$, $T^{n+1}x<t$ if and only if $T^nx\in\,]\frac1{k+t},\frac1k]$ for some $k$ (Problem~\ref{pb:gaussmeasure}(a)),
which gives the formula for $m_{n+1}$. Suppose $f_n$ is continuous and bounded by $M_n$ on $[0,1]$. The series
$\sum_k\frac{f_n(1/(k+t))}{(k+t)^2}$ of derivatives of the terms converges uniformly (it is bounded by $M_n\sum k^{-2}$), so $m_{n+1}$ is
differentiable with the stated derivative, continuous and bounded by $M_n\zeta(2)$. The same argument, one derivative higher,
shows by induction that $f_n$ is $C^1$, since $f_0=1$.

(b) $f_n(u_k)=\phi_n(u_k)/(1+u_k)$ and $1+u_k=\frac{k+t+1}{k+t}$, so
\[
\frac{(1+t)f_n(u_k)}{(k+t)^2}=\frac{1+t}{(k+t)(k+t+1)}\phi_n(u_k)=P_k(t)\phi_n(u_k).
\]
Also $\sum_kP_k(t)=(1+t)\sum_k\bigl(\frac1{k+t}-\frac1{k+t+1}\bigr)=1$.

(c) Let $\Lambda_n=\sup_{[0,1]}|\phi_n'|$. Differentiating (b) term by term (justified as in (a)),
\[
\phi_{n+1}'(t)=\sum_kP_k'(t)\bigl(\phi_n(u_k)-\phi_n(u_1)\bigr)-\sum_kP_k(t)\frac{\phi_n'(u_k)}{(k+t)^2},
\]
where we subtracted $\phi_n(u_1)\sum_kP_k'=0$. A computation gives
$P_k'(t)=\frac{k^2-k-(1+t)^2}{(k+t)^2(k+t+1)^2}$. For $k=1$ the first term vanishes. For $k\ge2$, $|k^2-k-(1+t)^2|\le k^2$, so
$|P_k'|\le\frac1{(k+1)^2}$, and $|\phi_n(u_k)-\phi_n(u_1)|\le\Lambda_n(u_1-u_k)=\Lambda_n\frac{k-1}{(1+t)(k+t)}\le\Lambda_n\frac{k-1}k$. For the second
sum, $\frac{P_k(t)}{(k+t)^2}=\frac{1+t}{(k+t)^3(k+t+1)}$ is at most $\frac12$ for $k=1$ (it equals $\frac1{(1+t)^2(2+t)}$) and at most
$\frac2{k^3(k+1)}$ for $k\ge2$. Hence $\Lambda_{n+1}\le c\Lambda_n$ with $c$ as stated. With $\frac1{k^3(k+1)}=\frac1{k^3}-\frac1{k^2}+\frac1k-\frac1{k+1}$
and $\frac{k-1}{k(k+1)^2}=\frac1{(k+1)^2}-\frac1{k(k+1)^2}$, $\frac1{k(k+1)^2}=\frac1k-\frac1{k+1}-\frac1{(k+1)^2}$, the sums are
$\sum_{k\ge2}\frac1{k^3(k+1)}=\zeta(3)-\zeta(2)+\frac12$ and $\sum_{k\ge2}\frac{k-1}{k(k+1)^2}=2\zeta(2)-3$, so $c=\frac12+2\zeta(3)-2\zeta(2)+1+2\zeta(2)-3=2\zeta(3)-\frac32=0.9041\ldots$

(d) $\phi_0=1+t$ has $\Lambda_0=1$, so $\Lambda_n\le c^n$ and $\max\phi_n-\min\phi_n\le c^n$. Since $\int_0^1f_n=1$,
$\int_0^1\phi_n(t)\frac{dt}{1+t}=1$: the number $\frac1{\log2}$ is the average of $\phi_n$ for the probability $\frac{dt}{(1+t)\log2}$,
so it lies between $\min\phi_n$ and $\max\phi_n$, and $|\phi_n(t)-\frac1{\log2}|\le c^n$. Divide by $1+t$, and integrate from $0$ to $t$:
$|m_n(t)-\log_2(1+t)|\le c^n\log(1+t)$.

(e) $P(a_{n+1}=k)=P(T^nx\in\,]\frac1{k+1},\frac1k])=m_n(\frac1k)-m_n(\frac1{k+1})$, which by (d) tends to
$\log_2(1+\frac1k)-\log_2(1+\frac1{k+1})$.

\tools Volume~1, Chapter~3 (uniform convergence, term-by-term differentiation); Volume~5, Chapter~12 (values of $\zeta$); Problem~\ref{pb:gaussmeasure}.
\leads Contraction arguments of this kind go back to Kuzmin (1928) and Lévy (1929); the constant $c$ here is far from the
truth, and a numerical calculation estimates the true rate by power iteration. Wirsing (1974) found the exact rate $0.3036630029\ldots$, the second eigenvalue of the transfer
operator on a suitable space. Babenko (1978) described the whole spectrum. Mayer later related this operator to the Selberg zeta
function of the modular group.
\end{solution}

\begin{solution}{pb:sqrtsum}
\idea Proposition~11.7 of Volume~5 is an identity between holomorphic functions, valid far to the left of $\operatorname{Re}s=1$. Read at $s=-\frac12$ it is exactly Euler--Maclaurin for $\sum\sqrt k$, with $\zeta(-\frac12)$
in the role of the constant.

\writeup Volume~5, Proposition~13.7 states that for $N,M\ge1$, $\operatorname{Re}s>1-2M$, $s\ne1$,
\[
\zeta(s)=\sum_{k\le N}k^{-s}+\frac{N^{1-s}}{s-1}-\frac{N^{-s}}2+\sum_{j=1}^M\frac{B_{2j}}{(2j)!}(s)_{2j-1}N^{1-s-2j}-\frac{(s)_{2M}}{(2M)!}\int_N^\infty\frac{\bar B_{2M}(x)}{x^{s+2M}}\dd x,
\]
with $(s)_i=s(s+1)\cdots(s+i-1)$ and $\bar B_{2M}$ bounded. Take $M=2$. For $s=-\frac12$: $\frac{N^{3/2}}{-3/2}=-\frac23N^{3/2}$, $-\frac{N^{1/2}}2$, and the $j=1$ term is $\frac1{12}\cdot(-\frac12)N^{-1/2}=-\frac1{24\sqrt N}$; the $j=2$ term
and the integral are $O(N^{-5/2})$. So $\sum_{k\le N}\sqrt k=\zeta(-\frac12)+\frac23N^{3/2}+\frac12\sqrt N+\frac1{24\sqrt N}+O(N^{-5/2})$. For $s=\frac12$: $\frac{N^{1/2}}{-1/2}=-2\sqrt N$, $-\frac1{2\sqrt N}$, the $j=1$ term is
$\frac1{12}\cdot\frac12N^{-3/2}=\frac1{24N^{3/2}}$, and the rest is $O(N^{-7/2})$. So $\sum_{k\le N}\frac1{\sqrt k}=\zeta(\frac12)+2\sqrt N+\frac1{2\sqrt N}-\frac1{24N^{3/2}}+O(N^{-7/2})$.

\tools{Volume~5, Theorem~13.5 (Euler--Maclaurin) and Proposition~13.7 (continuation of $\zeta$).}
\leads $\zeta(\frac12)<0$ although every term is positive: the constant is what is left after removing $2\sqrt N$, not a sum. By the functional equation (Volume~5, Chapter~12),
$\zeta(-\frac12)=-\frac{\zeta(3/2)}{4\pi}$. Ramanujan called such constants ``the constant of a series'', and Hardy's \emph{Divergent Series} (1949) makes the idea systematic.
\end{solution}

\section*{Sources and further reading}

\begin{itemize}
\item B.~Riemann, \emph{\"Uber die Anzahl der Primzahlen unter einer gegebenen Gr\"osse}, Monatsberichte der Berliner
Akademie, November 1859. Translated in the appendix of Edwards's book below.
\item H.~M.~Edwards, \emph{Riemann's Zeta Function}, Academic Press, 1974 (Dover reprint 2001). A line-by-line reading of
Riemann's paper; Chapter~1 contains both of his proofs of the functional equation.
\item T.~M.~Apostol, \emph{Introduction to Analytic Number Theory}, Springer, 1976, Chapters~11 and~12 (Dirichlet series,
Euler products, the Hurwitz zeta function).
\item E.~C.~Titchmarsh, \emph{The Theory of the Riemann Zeta-Function}, 2nd ed., revised by D.~R.~Heath-Brown, Oxford, 1986.
Chapter~2 gives seven proofs of the functional equation.
\item T.~Tao, \emph{The Euler--Maclaurin formula, Bernoulli numbers, the zeta function, and real-variable analytic continuation},
blog post, 2010. The smoothed sums of Proposition~\ref{prop:smoothed}, for every power $n^k$.
\item D.~J.~Platt and T.~S.~Trudgian, \emph{The Riemann hypothesis is true up to $3\cdot10^{12}$}, Bulletin of the London Mathematical
Society 53 (2021), 792--797.
\end{itemize}
